\documentclass[11pt]{article}
\usepackage[margin=1in]{geometry}
\usepackage{amsmath,amssymb}
\usepackage{booktabs}
\usepackage{hyperref}
\usepackage{xcolor}
\usepackage{longtable}

\title{Replication Report: Patra et al.\ (2022) --- PA1Pol$\beta\Delta$ Radiosensitivity}
\author{Ollie (LUCID replication subagent) \\ Backfilled by Kukla (subagent), 2026-07-06}
\date{Original triage: 2026-05-30 \quad PARTIAL-promotion pass: 2026-06-25}

\begin{document}
\maketitle

\begin{abstract}
We attempt a computational replication of Patra et al.\ 2022 (\textit{Radiat Oncol J} 40(1):66--78, DOI 10.3857/roj.2021.00689), which reports that PA1 ovarian cells carrying a truncated DNA polymerase $\beta$ (Pol$\beta\Delta$) are markedly more sensitive to $\gamma$-radiation than parental PA1. The paper's central radiobiological claim --- a $\sim$1.8-fold dose-modifying factor at 10~Gy driven by loss of base-excision-repair (BER) competence --- is qualitatively reproducible by LQ refit of the digitized Fig.~2. The mechanistic and computational half of the paper does not survive a spot-check: the deposited $\Delta$Pol$\beta$ cDNA cannot be translated into the protein the authors describe (frame-shifted, 6+ premature stops), deletion coordinates disagree across sections (208--301 / 208--304 / 211--339), docking scores are mislabelled as kcal/mol, and one input PDB (1WSR, aminomethyltransferase) is not a BER protein at all. \textbf{Verdict: PARTIAL} (coverage 6/10; agreement 6/10; paper-internal consistency 3/10).
\end{abstract}

\section{Paper's Headline Claim}
PA1Pol$\beta\Delta$ cells (isogenic PA1 line stably transfected with a domain-truncated Pol$\beta$ construct) show reduced clonogenic survival, elevated ROS and apoptosis, and altered cell-cycle distribution after $\gamma$-irradiation. Mechanism: loss of BER competence, supported by SWISS-MODEL structural model of Pol$\beta\Delta$ plus ClusPro/HDOCK dockings against 9 BER partners and against ss/dsDNA.

\section{What Was Exercised (Computational Portion Only)}
\begin{itemize}
  \item \textbf{Sequence check} (\texttt{01\_sequence\_check.py}): Biopython translation of WT and $\Delta$Pol$\beta$ cDNA from Suppl.\ Table~S1; alignment; frame check.
  \item \textbf{LQ refit} (\texttt{02\_lq\_fit.py}): weighted least-squares fit of $\mathrm{SF}(D) = \exp(-\alpha D - \beta D^2)$ to digitized Fig.~2 colony-forming data (two independent digitization passes).
  \item \textbf{Quantitative audit} (\texttt{03\_quantitative\_audit.py}): internal-consistency checks of tabulated ROS, cell-cycle, Annexin~V data and docking scores.
  \item \textbf{PDB structural audit} (\texttt{04\_pdb\_structural\_check.py}): downloaded all 11 cited PDB inputs (1TV9 + 10 BER partners), parsed with \texttt{Bio.PDB}, verified 7 canonical Pol$\beta$ active-site residues in 1TV9, cross-checked each partner's header vs.\ paper label.
\end{itemize}

\section{What Was NOT Exercised (and Why)}
\begin{itemize}
  \item \textbf{Wet-lab clonogenic survival / Western / DAPI / AO-PI / DCFDA / cell-cycle / Annexin~V.} No PA1 cells; out of scope for a LUCID computational triage.
  \item \textbf{Re-running ClusPro} (9 partners) --- public server AUP forbids parallel submissions; 1--12~h queue each; would produce the same arbitrary-unit output already published.
  \item \textbf{Re-running HDOCK} (4 protein--DNA dockings) --- same rate-limit issue.
  \item \textbf{Rebuilding SWISS-MODEL $\Delta$Pol$\beta$} --- pointless until the cDNA inconsistency is resolved; the deposited sequence encodes a 198-aa frame-shifted pseudo-protein, not a 238-aa domain-truncated polymerase.
  \item \textbf{Refitting raw colony counts} --- $n=3$ raw data not deposited.
  \item \textbf{Contacting authors} --- hard gate; not done.
\end{itemize}

\section{Reproduced Numbers}
\begin{table}[h]
\centering
\small
\begin{tabular}{lrr}
\toprule
Parameter & PA1 (WT) & PA1Pol$\beta\Delta$ \\
\midrule
$\alpha$ (Gy$^{-1}$)  & $0.045 \pm 0.025$ & $0.009 \pm 0.089$ \\
$\beta$ (Gy$^{-2}$)   & $0.0047 \pm 0.0021$ & $0.0284 \pm 0.0106$ \\
$\alpha/\beta$ (Gy)   & 9.6 & 0.30 \\
$D_{10}$ (Gy)         & \textbf{17.8} & \textbf{8.85} \\
SF$_2$                & 0.90 & 0.88 \\
\midrule
DMF at 10~Gy (vs.\ WT) & --- & \textbf{1.80} \\
\bottomrule
\end{tabular}
\caption{Our LQ refit of digitized Fig.~2. Qualitative claim of the paper (mutant more radiosensitive; 10~Gy selectively cytotoxic) is confirmed.}
\end{table}

\section{Critique --- Honest Assessment}
\subsection{Reproducibility Blockers}
\begin{enumerate}
  \item \textbf{Broken $\Delta$Pol$\beta$ cDNA.} The 595-nt sequence in Suppl.\ Table~S1 is a 413-nt deletion between WT codons 121 and 257, \emph{with a frameshift}. It encodes a 198-aa protein with $\geq 6$ internal stops, \emph{not} an in-frame 97-aa deletion at residues 208--304 as claimed. The SWISS-MODEL structural model of $\Delta$Pol$\beta$ therefore \emph{cannot} have been built from this cDNA. Every downstream mechanistic claim is built on an unverifiable input.
  \item \textbf{Three inconsistent deletion ranges.} Methods says 208--301; Results says 208--304; Discussion says 211--339 (which exceeds the full 335-aa WT length). Choose one.
  \item \textbf{PDB 1WSR is aminomethyltransferase, not a BER protein.} Silently mixes a mitochondrial folate-pathway enzyme into a 9-protein BER docking panel. Text mentions DNA ligase III in the pathway; there is no DNA-ligase PDB in the docking table.
  \item \textbf{Unit error.} ClusPro (weighted cluster score, AU) and HDOCK (native-like score, AU) outputs reported as kcal/mol. Not physical binding energies.
  \item \textbf{Undiscussed ROS anomalies.} $\Delta$Pol$\beta$ at 10~Gy has \emph{lower} ROS than at 5~Gy (dose-response inversion); PA1 has \emph{more} ROS at 10~Gy than the mutant (opposite of the BER-failure prediction).
  \item \textbf{Raw colony counts not deposited.} Trivially small ($\leq$30 numbers) but not shared, forcing digitization.
\end{enumerate}

\subsection{What Did Independently Verify}
Not everything is negative. We independently confirmed:
\begin{itemize}
  \item PDB 1TV9 is correctly identified as human Pol$\beta$ (title, chain length, resolution within rounding).
  \item All 7 canonical Pol$\beta$ active-site residues (D190, D192, D256, Y271, F272, N279, R283) present at correct positions in 1TV9 chain A coordinates.
  \item 9 of 10 cited docking-partner PDBs match their paper labels; only 1WSR is wrong.
  \item WT Pol$\beta$ cDNA translates cleanly to 335 aa; 22 spot-checked canonical residues all match.
  \item LQ refit qualitatively confirms $\sim$1.8$\times$ DMF at 10~Gy.
\end{itemize}

\section{Why PARTIAL (Not REPLICATED, Not NO-GO)}
The paper is a hybrid wet-lab + computational study. The wet-lab core (7 cell-biology figures) is gated: no public raw data, no cell lines, LUCID scope is computational. The computational periphery (sequence, LQ fit, structural inputs) \emph{is} reproducible and we did it. On the reproducible slice: radiobiology confirms; molecular mechanism disintegrates. PARTIAL is the honest verdict --- a REPLICATED label would overstate what the sequence/PDB audit found; NO-GO would understate the fact that the headline dose-modifying-factor claim actually holds up. This is the canonical LUCID wet-lab-gated pattern.

\section{Verdict}
\textbf{PARTIAL} --- radiobiology reproducible; mechanism non-reproducible at input layer. Coverage 6/10; agreement 6/10; paper-internal consistency 3/10.

\end{document}
